\section{Preprocessing: Which Rows and Quantities Were Kept, Dropped or Transformed}
\label{sec:preprocessing}

Both datasets are prepared the same way: the inputs are standardized on the training folds, the charge target is $\log_{10}(\Dch/M)$, the mass target is $M$ in solar masses, and whole curves are split into a frozen test set and GroupKFold folds. \Cref{tab:decisions} lists each decision of both datasets with its grounds from Sections~\ref{sec:eda-ns} and~\ref{sec:eda-bh}; the three choices a table row cannot carry are argued first.

\paragraph{The charge keeps every value: no floor.}
The charge target is $\log_{10}(\Dch/M)$, and every prepared charge is positive -- the smallest neutron-star charge is $\nsEdaDMin$, the smallest black-hole charge $\bhEdaDMin$ -- so the logarithm needs no floor, which would flatten the smooth, low-density neutron-star charges into a constant. On the logarithmic scale every charge carries the same relative weight across the $\nsEdaDOrders$ orders of magnitude $\Dch$ spans, as a relative-error metric asks (Section~\ref{sec:hypotheses}).

\paragraph{Whole curves are split: a frozen test set, then GroupKFold folds.}
Every split assigns whole curves. First, $\prepNsTestPercent\%$ of the curves are frozen as the test set, chosen with the fixed seed $\prepSeed$: $\prepNsTestCurves$ of the $\prepNsCurves$ neutron-star curves and $\prepBhTestCurves$ of the $\prepBhCurves$ black-hole curves. The test set is used once, for the final evaluation. The remaining curves are shuffled with the same seed into $\prepNsFolds$ GroupKFold folds for model selection, $\prepNsFoldMin$ to $\prepNsFoldMax$ neutron-star curves and $\prepBhFoldMin$ to $\prepBhFoldMax$ black-hole curves per fold. The standardization statistics are fitted on the training folds only, so no information from a held-out curve reaches the model.

\paragraph{Extrapolation is measured on held-out edge curves.}
The final model is trained on all curves outside the test set, so it covers the whole sampled domain. How far it carries beyond the sampled couplings is measured apart: every tuned candidate is refitted at its tuned settings without the edge curves of the coupling grid ($\prepNsAblationCurves$ neutron-star curves on the edges of the $(\beta, \lambda)$ grid and along its empty corner; the $\prepBhAblationCurves$ outermost black-hole curves) and without the test curves, and is scored on the edge curves it never saw. A map answers one step of the coupling grid beyond the sampled couplings, with a warning, when its edge curves keep at the 95th percentile at most one figure less than the interior test curves of that same refit and its worst edge row keeps a positive figure; otherwise its margin is zero and it refuses there. The margins and edge figures are in Section~\ref{sec:validity} and \cref{tab:zones}. The data foreshadowed the outcome: for neutron stars, predicting the rim from the interior costs about six times the error of predicting an interior curve ($\nsEdaSplitRimDM$ against $\nsEdaSplitCurvesDM$), while for black holes the outermost curves cost no more than any other ($\bhEdaSplitOuterCurvesLD$ against $\bhEdaSplitCurvesLD$).

% The decisions table's layout is paper.toml [tables] decisions (W-124): wide or narrow.
\input{93_preprocessing_tab_decisions_ns_\prepDecisionsLayout}
\FloatBarrier
